\section{Small-radius and large-outer-index expansions}
\label{ado:sec:asymptotics}
The global count removes an ambiguity from local expansions: the zeros
found near the first Fourier-mode zeros are the entire list of angular
zeros, not just selected branches.

\begin{theorem}[Two-term radial expansion]
\label{ado:thm:radialexpansion}
Fix an index satisfying \eqref{ado:eq:domain} with $d\ge2$, and put
\[
 p=\frac{c_{\svec}(d+1)}{C_{\svec}},\qquad
 q=\frac{c_{\svec}(d+2)}{C_{\svec}},\qquad
 \alpha_k=\frac{k\pi}{d}.
\]
The $k$th angular branch extends real analytically to $\rho=0$, and
\begin{align}
 \theta_k(\rho)={}&\alpha_k
 -\frac{p}{d}\sin\alpha_k\,\rho \notag\\
 &+\left(\frac{d+1}{d^2}p^2-\frac{2q}{d}\right)
       \sin\alpha_k\cos\alpha_k\,\rho^2
 +O_{\svec}(\rho^3).
 \label{ado:eq:radialexpansion}
\end{align}
In particular, every branch initially moves strictly toward smaller
angles as the radius increases from zero.
\end{theorem}
\begin{proof}
Divide the imaginary part of the disk series by $C_{\svec}\rho^d$.
The resulting analytic equation is
\[
 0=\sin(d\theta)+p\rho\sin((d+1)\theta)
                   +q\rho^2\sin((d+2)\theta)+O(\rho^3).
\]
At $(\rho,\theta)=(0,\alpha_k)$ its $\theta$ derivative is
$d(-1)^k\ne0$, so the analytic implicit-function theorem applies.
There are $d-1$ such branches; for small positive radius they exhaust
the global list from Theorem~\ref{ado:thm:angular}.

Write $\theta=\alpha_k+A\rho+B\rho^2+O(\rho^3)$ and divide by
$(-1)^k$. The coefficients of $\rho$ and $\rho^2$ are respectively
\[
 dA+p\sin\alpha_k,
 \qquad dB+(d+1)pA\cos\alpha_k+q\sin(2\alpha_k).
\]
Solving gives \eqref{ado:eq:radialexpansion}. The first derivative is
negative because $p>0$ and $0<\alpha_k<\pi$.
\end{proof}
Higher Taylor coefficients are obtained by the same finite coefficient
recursion. At each order the new unknown occurs with coefficient
$d(-1)^k$, while all remaining terms depend only on earlier
coefficients and finitely many moments. This is an exact algebraic
recurrence, not a proof of any global radial sign.

\subsection{Uniformity as the first index increases}
Fix $d\ge2$ and positive integer middle indices
$s_2,\ldots,s_{d-1}$. Let the first index be $a\in\bbN$ and the
last be $b>0$. Define
\[
 T_N(b)=\sum_{N\ge n_2>\cdots>n_d\ge1}
       \frac1{n_2^{s_2}\cdots n_{d-1}^{s_{d-1}}n_d^b},
 \qquad T_*=T_{d-1}(b)>0.
\]
The latter number is a single-chain value and is independent of $b$.
At depth two, $T_N(b)=H_N^{(b)}$ and $T_*=1$.

\begin{theorem}[First displacement, uniformly in the final order]
\label{ado:thm:largeouter}
Set $\lambda=d/(d+1)$ and $\mu=d/(d+2)$. As $a\to\infty$ through
positive integers, uniformly for $b>0$, $0\le\rho\le1$, and
$1\le k\le d-1$,
\begin{equation}\label{ado:eq:largeouter}
 \theta_k(\rho)
 =\alpha_k-\frac{T_d(b)}{dT_*}\sin\alpha_k\,
                      \rho\lambda^a
       +O\bigl(\rho^2\mu^a\bigr).
\end{equation}
At $\rho=0$, the branch is interpreted by its analytic extension.
The implied constant and the starting threshold for $a$ may depend on
the fixed depth and middle indices, but not on $b$, $\rho$, or $k$.
\end{theorem}
\begin{proof}
The normalized imaginary part is
\begin{equation}\label{ado:eq:outereq}
 \sin(d\theta)+\frac{T_d(b)}{T_*}\rho\lambda^a
                   \sin((d+1)\theta)+R(\theta),
\end{equation}
where
\[
 R(\theta)=\sum_{n\ge d+2}\frac{T_{n-1}(b)}{T_*}
        \rho^{n-d}\left(\frac dn\right)^a\sin(n\theta).
\]
All inner denominators are at least one, so
$0\le T_{n-1}(b)\le n^{d-1}$ uniformly for $b>0$. Choose the fixed
integer $a_0=d+3$. For $a\ge a_0$, termwise differentiation is
absolutely justified and
\begin{equation}\label{ado:eq:outertail}
 |R(\theta)|+|R'(\theta)|\le K\rho^2\mu^a,
 \quad
 K=\frac2{T_*}\sum_{n\ge d+2}n^d
                        \left(\frac{d+2}{n}\right)^{a_0}<\infty.
\end{equation}
Also $T_d(b)/T_*\le d^{d-1}/T_*$, independently of $b$.

On disjoint fixed small neighborhoods of the $\alpha_k$, the derivative
of the first term of \eqref{ado:eq:outereq} is bounded away from zero.
Its perturbation tends to zero uniformly by \eqref{ado:eq:outertail}.
Evaluating at $\alpha_k\pm K_1\rho\lambda^a$, for a sufficiently
large fixed $K_1$, gives opposite signs for all sufficiently large $a$.
This follows from the linear term $d(\theta-\alpha_k)$ and the
uniform bounds on the perturbations; $\mu/\lambda<1$. Thus each
neighborhood contains a unique root with displacement
$\delta=O(\rho\lambda^a)$. The global count identifies these roots
as all the branches in their established order.

Expand \eqref{ado:eq:outereq} at $\alpha_k$ and divide by $(-1)^k$.
The root equation becomes
\[
 0=d\delta+\frac{T_d(b)}{T_*}\rho\lambda^a\sin\alpha_k
     +O\bigl(\delta^2+\rho\lambda^a|\delta|
                         +\rho^2\mu^a\bigr).
\]
Because $\lambda^2<\mu$, the previously obtained displacement bound
makes the error $O(\rho^2\mu^a)$. Division by $d$ proves
\eqref{ado:eq:largeouter}. All constants used above are independent
of the final order.
\end{proof}
At depth two and unit radius this recovers the leading scale
$(2/3)^a$ already present in the manuscript's more detailed depth-two
expansion~\cite{repoSigned}. The added content is the simultaneous
all-depth statement, the exact branch indexing, and uniformity over the
entire range $b>0$ and $0\le\rho\le1$.
